\chapter{Introduction}
\label{ch:introduction}

\section{Topology control}
\label{sec:intro-topology-control}

Although a transmission grid is often represented as a fixed graph, operators
can change its topology. Inside each substation, transmission lines, generators,
and loads are connected to busbars. Moving an element from one
busbar to another, or splitting a substation into two separate electrical nodes,
changes the paths available to the power flow. The loading of the lines changes
as a result. This type of reconfiguration is called topology control.
\\
\\
This flexibility is attractive because it can reduce congestion without adding
new infrastructure. The usual alternatives are to redispatch generation,
curtail renewable production, or build new transmission lines. Redispatch and
curtailment cost money whenever they are used, while new lines are expensive
and take years to build. Topology control instead uses switching equipment that
is already installed, so the marginal cost of an action is close to
zero~\cite{fisher2008ots}. This option becomes even more useful as renewable
production increases, since variable generation makes power flows less
predictable and pushes existing lines closer to their limits.
\\
\\
Its low operating cost does not make topology control easy to use. In practice,
operators rely on a limited set of remedial actions prepared from experience and
offline studies. Searching every possible configuration during operation is not
realistic, which means that much of the available flexibility remains
unused~\cite{viebahn2022ai}.

\section{Why topology control should be automated}
\label{sec:intro-automation}

Two features of the problem make manual search difficult. The first is the
number of possible actions. If a substation has $d$ connected elements and two
busbars, it already has $2^{d}$ possible assignments. Considering several
substations together makes this number grow very quickly. Even the IEEE 14-bus
system used in this thesis has $178$ distinct unitary topology actions, while
the 36-substation WCCI grid in Chapter~\ref{ch:wcci-transfer} has more than
$66\,000$. Its action space must therefore be reduced before training
(Section~\ref{sec:wcci-action-space-reduction}).
\\
\\
The second difficulty is that these decisions are linked over time. An action
taken now changes the grid state and may affect which contingencies can be
handled later. Cooldown rules can also prevent another action at the same
substation for several time steps. The operator must therefore choose quickly
without considering only the immediate effect of an action.
\\
\\
Together, these properties make reinforcement learning (RL) a reasonable method
to investigate. An RL agent learns through repeated interaction with a
simulator, where each decision changes the next state and produces feedback.
There is no simple formula for the best action, but the objective is clear: keep
the grid operating and avoid a blackout. RTE created the Learning to Run a
Power Network (L2RPN) competitions to study this problem. Grid2Op supports these
experiments by providing the simulator, the load and generation time series,
and the operating rules~\cite{donnot2020grid2op,marot2021l2rpn}.

\section{Why use decentralized agents}
\label{sec:intro-decentralized}

Most published methods use one agent to control the full grid, but this is not
how large power systems are usually operated. The network is divided into
regions, and each control centre has detailed information about its own area
and a less complete view of the rest of the system~\cite{viebahn2022ai}. A
decentralized model follows this organization more closely: each actor receives
local observations, controls the substations in its region, and can operate
without a central controller at execution time.
\\
\\
The same organization also helps with scalability. A central actor must process
the full grid and produce scores for all available actions, so both its input
and output grow with the network. In a multi-agent system, each actor chooses
only among the actions of its own region. Adding another region therefore does
not require every actor to represent the full joint action space.
\\
\\
Decentralization does not remove the difficulty of the problem; it changes its
nature. Because all agents update their policies during training, the
environment appears non-stationary from the point of view of any
agent~\cite{tan1993independent}. Local observations also hide part of the grid,
so two actions that appear safe separately may be harmful when applied together.
\\
\\
Centralized training with decentralized execution helps address these issues.
During training, a critic can use the global state to evaluate the joint
behaviour of the agents, while each actor sees only the information that will be
available in its region at deployment time
~\cite{lowe2017maddpg,yu2022mappo}. The critic is removed after training, leaving
actors that operate independently. This is the setting used throughout the
thesis.

\section{Reinforcement learning for topology control}
\label{sec:intro-related-work}

Previous RL work on topology control can be grouped into four main approaches.
A broader review of graph-based RL for power systems is provided
by~\textcite{hassouna2024survey}.

\paragraph{Single-agent policies.}
The earliest approaches train one agent to control the whole grid.
\textcite{subramanian2021exploring} showed that deep RL with a selected action
set can outperform a do-nothing policy on L2RPN benchmarks. SMAAC, the winner
of the L2RPN WCCI 2020 challenge, combines actor--critic learning over
\emph{afterstates} with a hierarchical decision about where and how to
act~\cite{yoon2021smaac}. These methods can perform well, but the action space
must first be reduced, selected, or factorized.

\paragraph{Policies combined with heuristics.}
A second approach uses simple rules around a learned policy. A common choice is
to query the agent only when the maximum line loading $\rho_{\max}$ becomes
high, and to do nothing while the grid remains safe. PowRL combined this rule
with an overload manager and a reduced action set to win the L2RPN NeurIPS 2020
robustness track~\cite{chauhan2023powrl}. % \textcite{lehna2023comparative} found
% similar survival for advanced rule-based and RL agents, although RL was cheaper
% at execution time.

\paragraph{Hierarchical and centrally coordinated multi-agent methods.}
A third family of methods divides the decision between several components while
keeping a central coordinator. \textcite{manczak2023hrl} separate the decisions
of when, where, and how to act. \textcite{vandersar2023hmarl} use substation
agents under a higher-level selector, while \textcite{demol2025ccma} use a
central agent to choose between regional proposals. These methods reduce the
size of each decision, but a component still observes the full grid and
coordinates the agents. They therefore do not show whether local policies can
operate on their own.

\paragraph{Decentralized multi-agent control.}
Fully decentralized control has received less attention. Here, agents use local
observations and act without communication or a central selector, which is the
setting studied in this thesis. MARL2Grid addresses this gap by proposing a 
decentralized framework, because
earlier research mainly considered single-agent control and did not represent
the decentralized organization of grid operation
~\cite{marchesini2026marl2grid}.

\paragraph{Graph representations.}
Regardless of the control structure, the representation of the grid matters. A
flat vector hides the network structure and has a size tied to one grid. A graph
neural network (GNN) follows the grid connections and can use the same
parameters on networks of different sizes
~\cite{kipf2017gcn,velickovic2018gat,xu2018how}. However, a homogeneous graph may
omit connections that a topology action could create. Heterogeneous graphs can
reduce this \emph{busbar information asymmetry}~\cite{dejong2025gnntopo}.
Chapters~\ref{ch:methods} and~\ref{ch:experiments} compare these choices.

\section{Starting point: MARL2Grid and MAPPO}
\label{sec:intro-marl2grid}

This work builds on MARL2Grid~\cite{marchesini2026marl2grid}, the multi-agent
extension of RL2Grid~\cite{marchesini2025rl2grid}. MARL2Grid connects Grid2Op to
a multi-agent training pipeline and defines the regions, observations, action
spaces, rewards, and training loop used as the starting point of this thesis.

\paragraph{Grid partition.}
The grid is divided in advance into fixed, non-overlapping regions, with one
agent assigned to each region. In the IEEE 14-bus environment
(\texttt{bus14}), the three agents control substations $\{0,1,2,4\}$,
$\{3,6,7,8\}$, and $\{5,9,\ldots,13\}$. These regions contain $52$, $50$, and
$76$ non-idle actions, which together form the $178$ unitary actions of the full
grid. Each agent observes its own region and chooses an action at one of its
substations. The agents then act at the same time without seeing each other's
choices, and action $0$ always means do nothing.

\paragraph{Reward and evaluation metric.}
The reward is mainly based on survival, with $+1$ given for each survived step
and an additional term that penalizes overloads. Training therefore encourages
the agents to keep the episode alive. For one chronic, survival is defined as
the number of survived time steps divided by the maximum episode length, and
mean survival is this value averaged over the evaluated chronics.

\paragraph{MAPPO.}
Multi-agent PPO (MAPPO)~\cite{yu2022mappo} extends
PPO~\cite{schulman2017ppo} to centralized training with decentralized
execution. Each actor maps its local observation to a probability distribution
over its local actions. During training, PPO limits the size of each policy
update so that the actors do not change too abruptly. A centralized value
function can use the global state to evaluate the agents and estimate the
advantage of their actions. This global information is never given to the
actors, which continue to use only their local observations. Once training is
complete, the centralized value function is no longer needed and deployment
remains decentralized.

\paragraph{Initial baseline.}
The initial model uses flat multilayer-perceptron actors and one critic trained
on \texttt{bus14}. Without further tuning, this baseline is unstable and has low
survival. It is also tied to one grid because the actor input size depends on
the number of substations, so its weights cannot be used directly on a larger
network. The next chapters address these limitations in sequence: they first
stabilize training, then reduce unnecessary interventions, and finally replace
the fixed-size actor with a graph-based model.

\section{Outline}
\label{sec:intro-outline}

\newrevision{The thesis is organized in two stages. Chapters~\ref{ch:baseline}
and~\ref{ch:sparse-control} first focus on learning a reliable decentralized
controller on \texttt{bus14}. Chapter~\ref{ch:baseline} establishes a stable
flat MAPPO baseline, and Chapter~\ref{ch:sparse-control} studies how to reduce
unnecessary interventions. Its preliminary WCCI stress test also reveals the
limitations of applying a fixed-size actor to a larger grid.}
\\
\\
{\color{red}Chapters~\ref{ch:methods}--\ref{ch:screen-f-transfer} then develop
and evaluate cross-grid transfer. Chapter~\ref{ch:methods} introduces the graph
actor and candidate-action scorer, and Chapter~\ref{ch:experiments} evaluates
their main design choices on \texttt{bus14}. Chapter~\ref{ch:encoder-inputs}
makes the physical inputs more comparable across grids, while
Chapter~\ref{ch:wcci-transfer} defines the WCCI environment and evaluation
protocol. Chapter~\ref{ch:encoder-transfer-bridge} tests whether the encoders
selected by Screens A--E remain useful when their grid-specific action heads
are replaced. Chapter~\ref{ch:screen-f-transfer} then studies the stricter case
in which the complete candidate-action policy is transferred, with and without
target adaptation. Chapter~\ref{ch:discussion-conclusion} closes the thesis by
answering the research questions and discussing the limits of the evidence.}
